% sections/appendices.tex -- Technical appendices: conventions, graphs, zero-dimensional sum rules, reproducibility,
% data and code
% \appendix is issued here (main.tex does not issue it).
\appendix

% =================================================================================================================
\section{Conventions and normalizations}
\label{app:conventions}
% A Markdown summary of this appendix is kept in ../doc/CONVENTIONS.md; change both together.

These conventions are also summarized, with the names used by the code, in the file \texttt{doc/CONVENTIONS.md} of
the code repository (Appendix~\ref{app:data}).

\paragraph{Couplings.}
The renormalized coupling $g$ is defined by $\Gamma^{(4)}_R(0)=m\,g$ [Eq.~\eqref{eq:normcond}], with the $O(N)$ tensor
$(\delta_{ij}\delta_{kl}+\delta_{ik}\delta_{jl}+\delta_{il}\delta_{jk})/3$ factored out. The integrals are expanded in
$u=g/8\pi$ and the RG functions are quoted in
\begin{gather}
\ubar=\frac{N+8}{6}\,u=\frac{N+8}{48\pi}\,g ,\nonumber\\
\beta(\ubar)\equiv m\frac{\partial\ubar}{\partial m}=\frac{N+8}{6}\,\beta_u ,
\label{eq:ubar}
\end{gather}
the normalized coupling of Refs.~\cite{BNGM1976,BNM1978}, which is the $\tilde u$ of Ref.~\cite{Phi4tools2024}. With
it $\beta(\ubar)=-\ubar+\ubar^2+O(\ubar^3)$ for every $N$. We write $\eta_{\phi^2}$ for the anomalous dimension of $\phi^2$ ($\eta_2$ in
Ref.~\cite{Ribeca2003thesis}), so that $1/\nu=2-\eta+\eta_{\phi^2}$, as in Ref.~\cite{Phi4tools2024}.

\paragraph{Integrals.}
Our integrals have one $d^3k/\pi^2$ per loop~\cite{NMB1977}, that is $8\pi$ times the measure $d^3k/(2\pi)^3$; with
this choice each loop comes with one power of $u$, the one-loop bubble at zero momentum equals one, and the low-order
coefficients are rational. (In Ref.~\cite{Ribeca2003thesis} ``Nickel's normalization'' refers to the coupling
$\ubar$ rather than to the measure; integrals normalized with $d^3k/(2\pi)^3$ per loop are $(8\pi)^{-L}$ times ours.)
The propagators of $B_L$, $H_L$ and $C_L$ have the renormalized mass, $m=1$, whereas those of
Refs.~\cite{NMB1977,Ribeca2003thesis,Phi4tools2024}, and of the intermediate-scheme integrals we use for comparison,
have the mass of the intermediate scheme, $m_I=1$. Individual integrals in the two conventions differ by powers of
$m^2/m_I^2$; the RG functions do not. Reference~\cite{Phi4tools2024} quotes graph values $V_G$ that include its weight
factor and a factor $(16\pi)^{L}$ per graph with $L$ loops, with the measure $d^3k/(2\pi)^3$. At $N=1$ the sum of
its intermediate-scheme four-point values at $L$ loops equals $(2^{L}/3)\sum_G S_G I_G$, where $I_G$ is the integral of
$G$ in our normalization, in the same scheme and with $m_I=1$; this is the relation behind the comparisons of
intermediate-scheme order sums in Sec.~\ref{sec:results-validation}.

\paragraph{Weights.}
The weight $S_G=T_G(N)\,E!/|\mathrm{Aut}\,G|$ of Sec.~\ref{sec:methods-graphs} includes the $O(N)$ factor $T_G(N)$,
normalized to $T_G(1)=1$ as are the symmetry factors of Ref.~\cite{Phi4tools2024}. The insertion of $\tfrac12\phi^2$ is
normalized by $\Gamma^{(2,1)}_R(0;0)=1$, and $\partial\Gamma^{(a,b)}/\partial m_0^2=\Gamma^{(a,b+1)}$ at zero insertion
momentum, so that, e.g., $C_1>0$.

\paragraph{From the per-order sums to the series.}
Writing $[\ubar^k]F$ for the coefficient of $\ubar^k$ in $F$, Eqs.~\eqref{eq:csrel}, \eqref{eq:series}
and~\eqref{eq:ubar} give, at each order,
\begin{align}
[\ubar^{L+1}]\beta&=2(-1)^{L+1}\Bigl(\frac{6}{N+8}\Bigr)^{L}\bigl(B_L-2H_L\bigr)+\dots,\nonumber\\
[\ubar^{L}]\eta&=2(-1)^{L}\Bigl(\frac{6}{N+8}\Bigr)^{L}H_L+\dots,\nonumber\\
[\ubar^{L}]\eta_{\phi^2}&=2(-1)^{L}\Bigl(\frac{6}{N+8}\Bigr)^{L}C_L+\dots,\label{eq:linear}\\
[\ubar^{L}]\nu&=\frac{(-1)^{L+1}}{2}\Bigl(\frac{6}{N+8}\Bigr)^{L}\bigl(C_L-H_L\bigr)+\dots,\nonumber
\end{align}
where the dots stand for products of sums of lower order ($H_1=0$). The coefficient of $\ubar^{L+1}$ in $\beta$ thus
depends on $B_L$ and $H_L$, those of $\ubar^{L}$ in $\eta$ on $H_L$, and in $\nu$ on $C_L$ and $H_L$; the errors of the
new sums enter linearly, with the prefactors of Eq.~\eqref{eq:linear}. The reference values of $B_L$, $H_L$, $C_L$ used
in Sec.~\ref{sec:results} are obtained by solving these relations order by order for the sums, starting from the
series of Ref.~\cite{Phi4tools2024}~\check{S1}.

At the lowest orders the sums are known exactly,
\begin{gather}
B_1=\frac{N+8}{12},\qquad B_2=\frac{5N+22}{54},\qquad H_2=\frac{N+2}{243},\nonumber\\
C_1=\frac{N+2}{12},\qquad C_2=\frac{N+2}{36},
\label{eq:exactsums}
\end{gather}
that is $3/4$, $1/2$, $1/81$, $1/4$ and $1/12$ at $N=1$, and give
\begin{align}
\beta(\ubar)&=-\ubar+\ubar^2-\frac{4(41N+190)}{27(N+8)^2}\,\ubar^3+O(\ubar^4),\nonumber\\
\eta(\ubar)&=\frac{8(N+2)}{27(N+8)^2}\,\ubar^2+O(\ubar^3),\nonumber\\
\eta_{\phi^2}(\ubar)&=-\frac{N+2}{N+8}\,\ubar+\frac{2(N+2)}{(N+8)^2}\,\ubar^2+O(\ubar^3),\label{eq:lowseries}\\
\nu(\ubar)&=\frac12+\frac{N+2}{4(N+8)}\,\ubar\nonumber\\
&\quad+\frac{(N+2)(27N-38)}{216(N+8)^2}\,\ubar^2+O(\ubar^3),\nonumber
\end{align}
in agreement with Ref.~\cite{Phi4tools2024}. At $N=1$ the coefficients are $-308/729$ ($\ubar^3$ in $\beta$),
$8/729$ ($\ubar^2$ in $\eta$), $-1/3$ and $2/27$ ($\eta_{\phi^2}$), and $1/12$ and $-11/5832$ ($\nu$). In the notation of
Eq.~\eqref{eq:series-form}, $p_{\beta,3}=-4(41N+190)/27$, $p_{\eta,2}=8(N+2)/27$, $p_{\nu,1}=(N+2)/4$ and
$p_{\nu,2}=(N+2)(27N-38)/216$.

% =================================================================================================================
\section{Graphs}
\label{app:graphs}

\paragraph{Counts.}
Table~\ref{tab:graphcounts} gives, for every loop order, the numbers of 1PI graphs and of the integrands of
Eq.~\eqref{eq:BHC} that do not vanish identically. One-vertex-reducible (1VR)
graphs vanish for all three kinds, and so do, for $H_L$, the graphs with both external legs on one vertex
(Sec.~\ref{sec:methods-rbar})~\check{I2}; through seven loops no other graph vanishes, so that the numbers of integrands are the
numbers of the remaining graphs; at eight and nine loops the expansion confirmed that every remaining graph has a
non-vanishing integrand (no class was dropped at the build). At ten loops the table gives the numbers of remaining
graphs~\check{G5}. In the runs reported in Sec.~\ref{sec:results} each class is a single graph.

\begin{table*}
\caption{\label{tab:graphcounts}%
Numbers of 1PI graphs per loop order $L$ and of non-vanishing integrands. $|\mathcal G_4(L)|$, $|\mathcal G_2(L)|$:
graphs without self-loops, one-vertex-reducible ones included; 1VR: one-vertex-reducible graphs among them; one vertex:
two-point graphs with both legs on one vertex (generalized tadpoles), 1VR ones included; $B$, $H$, $C$: integrands of
Eq.~\eqref{eq:BHC} that survive the algebraic pruning, i.e.\ graphs that are not 1VR and, for $H$, do not have both
legs on one vertex; with self-loops: complete two-point lists including graphs with self-loops, used only for
checks. \check{G1}}
\begin{ruledtabular}
\begin{tabular}{rrrrrrrrrr}
$L$ & $|\mathcal G_4(L)|$ & 1VR & $B$ & $|\mathcal G_2(L)|$ & 1VR & one vertex & $H$ & $C$ & with self-loops \\
\hline
1 & 1        & 0    & 1                      & 0    & --   & --   & --                    & 1\footnotemark[1]     & 1    \\
2 & 2        & 1    & 1                      & 1    & 0    & 0    & 1                     & 1                     & 2    \\
3 & 8        & 2    & 6                      & 2    & 0    & 1    & 1                     & 2                     & 5    \\
4 & 27       & 8    & 19                     & 6    & 1    & 2    & 4                     & 5                     & 17   \\
5 & 129      & 33   & 96                     & 19   & 3    & 7    & 11                    & 16                    & 68   \\
6 & 660      & 158  & 502                    & 75   & 13   & 21   & 50                    & 62                    & 313  \\
7 & 3986     & 836  & 3150                   & 317  & 52   & 85   & 209                   & 265                   & 1632 \\
8 & 26\,540  & 4851 & 21\,689 & 1622 & 235  & 356  & 1157  & 1387  & 9685 \\
9 & 197\,022 & 30\,975 & 166\,047\footnotemark[2] & 9184 & 1146 & 1804 & 6776  & 8038  & 63\,774 \\
10 & 1\,593\,254 & 215\,576 & 1\,377\,678\footnotemark[2] & 59\,614 & 6217 & 10\,064 & 46\,102\footnotemark[2] & 53\,397\footnotemark[2] & 463\,933 \\
\end{tabular}
\end{ruledtabular}
\footnotetext[1]{The one-loop tadpole with two insertions.}
\footnotetext[2]{Graphs that are not 1VR (and, for $H$, do not have both legs on one vertex); all their integrands
survive the pruning, and every one of them was sampled~\check{G5}.}
\end{table*}

\paragraph{Generator.}
The generation is described in Sec.~\ref{sec:methods-generation}. At eight loops the
26\,540 four-point graphs have their legs distributed as $(2,2)$ (2254 graphs), $(2,1,1)$ (14\,614) and $(1,1,1,1)$
(9672); at nine loops the 9184 two-point graphs as $(2)$ (1804) and $(1,1)$ (7380)~\check{G1}.

\paragraph{Twins.}
Besides the seven-loop numbers of Sec.~\ref{sec:methods-construction}, the four-point lists at five and six loops contain
104 and 518 distinct unsubtracted integrals for 129 and 660 graphs, and the eight-loop two-point list 1512 for 1622.
Merging twins, after a pointwise check of their subtracted integrands, would save the corresponding samples.

\paragraph{Notation.}
A Nickel index~\cite{NMB1977} lists, for each vertex in turn, an \texttt{e} for every external leg and the labels of the
vertices with equal or higher label to which it is joined, once per line, followed by a bar; the canonical index is the
lexicographic minimum over the labellings of the vertices, with $|<\texttt{e}<0<1<\dots$. For example, the one-loop
four-point bubble is \texttt{ee11|ee|} and the two-loop sunset \texttt{e111|e|}. Our generator, like the one of
Ref.~\cite{Ribeca2003thesis} and the programs built on it~\cite{GuidaRibeca2006}, writes a graph instead as its number
of vertices followed by the run-length-coded upper triangle of its symmetric matrix of line multiplicities, with
optional sections for vertex colours (external legs) and self-loops. Those programs canonicalized it as a
lexicographically minimal adjacency matrix, and their labels such as D40L5N78 (a four-point graph with five loops,
line 78 of their list) are list positions, not graph invariants. For instance, the string \texttt{6(.2 1\^{}2.2 1\^{}6.1 1\^{}2)} of D40L5N78
is the graph \texttt{ee12|345|345|e4|5|e|}. A table that maps every historical label to its Nickel index accompanies
the graph lists (Appendix~\ref{app:data}).

% =================================================================================================================
\section{Zero-dimensional sum rules}
\label{app:sumrules}

In zero dimensions every Feynman integral with unit propagators equals one. The sum $\sum_G S_G$ of the weights over a
complete list is therefore the corresponding coefficient of the vertex functions of the zero-dimensional $O(N)$
model~\cite{CvitanovicLautrupPearson1978,Ribeca2003thesis,Borinsky2017}, with partition function
\begin{equation}
Z=\int d^N\phi\;\exp\Bigl[-\frac{m_0^2}{2}\,\phi^2-\frac{g}{24}\,(\phi^2)^2\Bigr].
\label{eq:zerod}
\end{equation}
With $Q=\langle\phi^2\rangle/N$ and $A=\langle(\phi^2)^2\rangle/[N(N+2)]$ one has $\Gamma^{(2)}_0=1/Q$ and
$\Gamma^{(4)}_0=-3(A-Q^2)/Q^4$, and, as formal power series in $g$,
\begin{equation}
\langle(\phi^2)^k\rangle=\frac{I_k}{I_0},\quad
I_k=\sum_{n\ge0}\frac{(-g/24)^n}{n!}\,\frac{P_{k+2n}(N)}{m_0^{2(k+2n)}},
\label{eq:zerodmoments}
\end{equation}
with $P_j(N)=N(N+2)\cdots(N+2j-2)$. For the lists without self-loops (one-vertex-reducible graphs included), the bare
mass is shifted, $m_0^2=1-(N+2)g/6$, so that every graph with a self-loop cancels against a graph with a mass
counterterm; then
\begin{gather}
\Gamma^{(2)}_0=1+\sum_{L\ge2}(-1)^{L+1}\Sigma_{2,L}\,g^L,\nonumber\\
\Gamma^{(4)}_0=\sum_{L\ge0}(-1)^{L}\Sigma_{4,L}\,g^{L+1},
\label{eq:sumrule}
\end{gather}
with $\Sigma_{E,L}=\sum_{G\in\mathcal G_E(L)}S_G$. With $m_0^2=1$ the same expansion gives the sums over all 1PI graphs,
self-loops included. Table~\ref{tab:sumrules} gives the polynomials $\Sigma_{E,L}(N)$. Through eight loops they are those
of the zero-dimensional computation of Ref.~\cite{Ribeca2003thesis}, whose table stops at $\Gamma^{(4)}$ with seven
loops and $\Gamma^{(2)}$ with seven, and whose original output extends to eight loops in both. The printed table has two
sign misprints, absent from the original output: the last four terms of the six-loop $\Gamma^{(2)}$ numerator and the
term in $N^3$ of the three-loop $\Gamma^{(2,1)}$ numerator have the same sign as the other terms. The nine- and
ten-loop rows were obtained here from Eqs.~\eqref{eq:zerod}--\eqref{eq:sumrule}, which reproduce every coefficient
of the original computation at lower orders.

The lists used here reproduce every row of Table~\ref{tab:sumrules} exactly, as polynomials in $N$, through ten
loops for both $\Gamma^{(4)}$ and $\Gamma^{(2)}$, up to the 1\,593\,254 ten-loop four-point graphs; the $O(N)$ factor of
every graph is counted exactly by carrying the partial pairings of its vertices from vertex to vertex, which agrees
graph by graph with the enumeration of all $3^V$ pairings wherever both run~\check{G2}. The lists with self-loops reproduce the bare series,
whose coefficients at $N=1$ are $1/2$, $5/12$, $5/6$, $115/48$, $625/72$, $10\,775/288$, $80\,275/432$,
$21\,615\,925/20\,736$, $67\,531\,025/10\,368$ and $620\,085\,875/13\,824$ at $L=1$--10. They also reproduce the sum rules of
$\Gamma^{(2,1)}$ of Ref.~\cite{Ribeca2003thesis}, obtained by inserting $\phi^2$ on any line of a graph without
self-loops and on the self-loop of a graph with exactly one, at five and six loops for all $N$, the seven-loop
polynomial of the original output, and its eight-loop value at $N=1$.

\begin{table*}
\caption{\label{tab:sumrules}%
Zero-dimensional sum rules: $\Sigma_{E,L}(N)=\sum_{G\in\mathcal G_E(L)}S_G$ over the $L$-loop 1PI graphs without self-loops,
one-vertex-reducible graphs included [Eq.~\eqref{eq:sumrule}], written as a numerator polynomial over a denominator
$D$, and its value at $N=1$. Rows to eight loops: Ref.~\cite{Ribeca2003thesis} (printed table and original output);
the nine- and ten-loop rows: this work, from the lists and from the exact zero-dimensional series, which agree.
\check{G2}}
\footnotesize
\begin{ruledtabular}
\begin{tabular}{llp{0.62\textwidth}rr}
$E$ & $L$ & Numerator & $D$ & $N=1$ \\
\hline
4 & 1 & $8+N$ & 6 & $3/2$ \\
  & 2 & $108+26N+N^2$ & 36 & $15/4$ \\
  & 3 & $2048+724N+62N^2+N^3$ & 216 & $105/8$ \\
  & 4 & $47648+21632N+2928N^2+124N^3+N^4$ & 1296 & $893/16$ \\
  & 5 & $1297664+702864N+127392N^2+9044N^3+220N^4+N^5$ & 7776 & $8795/32$ \\
  & 6 & $40362304+24950944N+5549104N^2+546408N^3+23404N^4+358N^5+N^6$ & 46656 & $97987/64$ \\
  & 7 & $1410726400+966932160N+250483648N^2+31004768N^3+1898240N^4+53312N^5+546N^6+N^7$ & 279936 & $10951025/1152$ \\
  & 8 & $54747705344+40784744192N+11897634304N^2+1746112576N^3+137234272N^4+5654416N^5+110208N^6+792N^7+N^8$ & 1679616 & $49985915/768$ \\
  & 9 & $2336870225920+1865178040064N+598669644160N^2+100305560704N^3+9494392864N^4+511500848N^5+14961264N^6+211032N^7+1104N^8+N^9$ & 10077696 & $2245562203/4608$ \\
  & 10 & $108856221154304+92103242480128N+31990963159296N^2+5962312144640N^3+652186499392N^4+42906530592N^5+1667820720N^6+36012864N^7+379788N^8+1490N^9+N^{10}$ & 60466176 & $4057808535/1024$ \\
\hline
2 & 2 & $2+N$ & 18 & $1/6$ \\
  & 3 & $10+7N+N^2$ & 54 & $1/3$ \\
  & 4 & $280+232N+52N^2+3N^3$ & 648 & $7/8$ \\
  & 5 & $2552+2384N+688N^2+71N^3+2N^4$ & 1944 & $211/72$ \\
  & 6 & $113888+116608N+39952N^2+5680N^3+320N^4+5N^5$ & 23328 & $3413/288$ \\
  & 7 & $1498912+1650832N+641936N^2+112904N^3+9260N^4+316N^5+3N^6$ & 69984 & $48323/864$ \\
  & 8 & $90828928+106190272N+45557472N^2+9353504N^3+982000N^4+51060N^5+1134N^6+7N^7$ & 839808 & $3123017/10368$ \\
  & 9 & $778012480+956252384N+444020448N^2+102678120N^3+12884816N^4+882730N^5+31028N^6+473N^7+2N^8$ & 1259712 & $9443467/5184$ \\
  & 10 & $118963074560+152597681664N+75647848448N^2+19237783296N^3+2767156512N^4+231239968N^5+10977808N^6+273456N^7+2982N^8+9N^9$ & 30233088 & $506798407/41472$ \\
\end{tabular}
\end{ruledtabular}
\end{table*}

% =================================================================================================================
\section{Reproducibility}
\label{app:reproducibility}

Every claim of this paper that a script can verify is marked in the text with the identifier of its check, as in
\check{G2}. The checks are part of the code repository (Appendix~\ref{app:data}), in its directory \texttt{checks/}: one
script per identifier, which states the claim as it is made here, its inputs in the repository and the expected result,
prints the numbers it obtains, and ends with \texttt{PASS} or \texttt{FAIL}. Table~\ref{tab:checks} lists them. The
command
\begin{verbatim}
make -C checks
\end{verbatim}
runs them all and writes a report, \texttt{checks/REPORT.md}; \texttt{make -C checks quick} runs those that take less
than about a minute, and \texttt{python3 checks/run.py G2} runs one. They need Python~3.6 or later with its standard
library, GNU make and a C++17 compiler with OpenMP, which build the CPU integrator and the compiled helpers on first
use. Where a check needs what the machine lacks, a
GPU for the live part of N2 or the \textsc{mpmath} library for I1, it reports \texttt{SKIP} for that part and verifies
instead the stored outputs of the runs that did it, which the repository keeps (\texttt{results/}).

The numbers that describe the pooled Monte Carlo data are not typed in the sources of the paper. The bodies of
Tables~\ref{tab:validation}, \ref{tab:cost}, \ref{tab:newsums}--\ref{tab:newcoeffnu}, \ref{tab:routes},
\ref{tab:exponents} and~\ref{tab:largeorder}, the rows of this work in Table~\ref{tab:effort}, Figs.~\ref{fig:cost}
and~\ref{fig:costhist}, and the values quoted from them in the text are written by \texttt{make -C series} from the
ledger of pooled per-class statistics (Sec.~\ref{sec:methods-ledger}) and from the stored runs of route B. Check S2
regenerates all of them except Table~\ref{tab:routes}, which check S4 regenerates, and both verify that they are
identical to those printed here. The exceptions are the numbers of individual early runs and tests in
Sec.~\ref{sec:results-validation} (the round-off example of $H_4$ and the test of fused multiply-add contraction),
which are historical; where a check exists, it verifies them from their stored outputs.

\begin{table*}
\caption{\label{tab:checks}%
The reproducibility checks of the repository (directory \texttt{checks/}). Where: the sections and appendices whose
claims cite the check. Needs: CPU, a computation on a CPU with the graph tools or the CPU build of the integrator; GPU,
a computation on a GPU; stored, only the stored outputs of the repository (the ledger and the reference data of the
checks).}
\begin{ruledtabular}
\begin{tabular}{lp{0.2\textwidth}p{0.55\textwidth}l}
Id & Where & Claim & Needs \\
\hline
\checkid{G1} & \ref{sec:methods-generation}, \ref{sec:methods-checks}, App.~\ref{app:graphs} &
  Every graph list holds distinct, connected 1PI graphs in canonical Nickel form, in the numbers of
  Tables~\ref{tab:generation} and~\ref{tab:graphcounts} & CPU \\
\checkid{G2} & \ref{sec:methods-checks}, \ref{sec:methods-validation}, App.~\ref{app:sumrules} &
  The lists satisfy the zero-dimensional sum rules exactly, as polynomials in $N$, through ten loops, and reproduce
  every row of Table~\ref{tab:sumrules} & CPU \\
\checkid{G3} & \ref{sec:methods-checks} &
  The lists equal those of Ref.~\cite{Phi4tools2024} as sets: $\Gamma^{(4)}$ to seven loops, $\Gamma^{(2)}$ to eight
  & CPU \\
\checkid{G4} & \ref{sec:methods-checks} &
  The weights and $O(N)$ polynomials equal those of Refs.~\cite{Ribeca2003thesis,GuidaRibeca2006} on all graphs with
  $L\le5$ and on the four-loop tables & CPU \\
\checkid{G5} & \ref{sec:methods-tropical}, App.~\ref{app:graphs} &
  The counts of Table~\ref{tab:graphcounts} follow from the lists (1VR by a cut-vertex test, checked against the
  lists' flags), and the ledger holds that many classes for every order with data & CPU \\
\checkid{I1} & \ref{sec:methods-checks}, \ref{sec:methods-validation} &
  The expanded forest formula, Eq.~\eqref{eq:integrand}, agrees pointwise with an independent evaluation in 100-digit
  arithmetic & CPU\footnotemark[1] \\
\checkid{I2} & \ref{sec:methods-construction}, \ref{sec:methods-checks}, \ref{sec:methods-validation},
  App.~\ref{app:graphs} &
  One-vertex-reducible graphs, graphs with self-loops and generalized tadpoles vanish after the pruning & CPU \\
\checkid{I3} & \ref{sec:methods-construction}, \ref{sec:methods-checks}, \ref{sec:methods-validation} &
  Twins (2-isomorphic graphs) have pointwise equal subtracted integrands & CPU \\
\checkid{I4} & \ref{sec:methods-construction} &
  A run that loads the built classes from the cache is bitwise identical to one that builds them & CPU \\
\checkid{N1} & \ref{sec:methods-rng} &
  Our Philox4x32-10 reproduces the known-answer vectors of Ref.~\cite{Salmon2011Random123} & CPU \\
\checkid{N2} & \ref{sec:methods-kernel}, \ref{sec:methods-validation}, \ref{sec:results-validation} &
  On identical random numbers the GPUs reproduce the per-class means of the CPU to round-off (golden runs)
  & CPU, GPU\footnotemark[2] \\
\checkid{N3} & \ref{sec:methods-checks}, \ref{sec:results-validation} &
  The exact two-loop sums $B_2=1/2$, $H_2=1/81$ and $C_2=1/12$ are reproduced within the errors & stored \\
\checkid{N4} & \ref{sec:results-validation} &
  The sums at three to seven loops agree within the errors with those derived from Ref.~\cite{Phi4tools2024}; the
  statistics of the pulls & stored \\
\checkid{N5} & \ref{sec:methods-ledger}, \ref{sec:results-validation} &
  The ledger is rebuilt bit for bit from the archive of its runs, and every run is consistent class by class with
  the pool of the runs before it (per-class $\chi^2$) & stored\footnotemark[3] \\
\checkid{N6} & \ref{sec:results-validation} &
  FP64 round-off does not bias the order sums: paired double and extended-precision evaluation of 55 classes (the 43 most exposed of $B_8$, $H_8$,
  $C_8$, $H_9$, $C_9$ and $C_{10}$, and 12 controls) finds no significant difference, and at most $4\times10^{-6}$ of an order's error summed over its classes & CPU \\
\checkid{N7} & \ref{sec:methods-tropical}, \ref{sec:results-validation} &
  Eight replicas of each of $B_6$, $H_6$, $C_6$, $B_7$, $H_7$ and $C_7$ scatter as their quoted errors say, order
  sums and class by class & stored \\
\checkid{N8} & \ref{sec:results-validation} &
  The round-off scales of the four classes of $H_4$ quoted as an example in Sec.~\ref{sec:results-validation}
  & stored \\
\checkid{S1} & \ref{sec:methods-validation}, \ref{sec:results}, App.~\ref{app:conventions} &
  The Callan--Symanzik relations map the sums derived from Ref.~\cite{Phi4tools2024} back to its coefficients, to
  $10^{-15}$ & stored \\
\checkid{S2} & \ref{sec:methods-ledger}, \ref{sec:resummation-criteria}, \ref{sec:resummation-errors},
  \ref{sec:methods-validation}, \ref{sec:results-validation}, \ref{sec:results-performance}, \ref{sec:results-new},
  \ref{sec:resummation-results} &
  Every number of the generated tables and values is regenerated from the ledger by \texttt{make -C series} & stored \\
\checkid{S3} & \ref{sec:methods-checks}, \ref{sec:methods-validation} &
  The large-order ratios of Table~\ref{tab:largeorder} follow from the coefficients and the constants of
  Sec.~\ref{sec:methods-validation} & stored \\
\checkid{S4} & \ref{sec:methods-validation}, \ref{sec:results-new} &
  Route B reproduces the coefficients of Ref.~\cite{Phi4tools2024} from its per-graph values, its eight-vertex
  sums agree with those of Ref.~\cite{Phi4tools2024} and its ninth- and tenth-order coefficients with route A
  (Table~\ref{tab:routes}) & stored \\
\checkid{S5} & \ref{sec:methods-validation}, \ref{sec:results-new} &
  The terms in $N^L$ of every order sum vanish class by class, as the limit $N\to\infty$ requires, and those in
  $N^{L-1}$ of every order of the ledger agree with the $1/N$ expansion, which reproduces the leading coefficients of
  Ref.~\cite{Phi4tools2024}; the polynomials $p_{X,k}(N)$ have the degrees and leading coefficients that follow
  & stored \\
\end{tabular}
\end{ruledtabular}
\footnotetext[1]{With the \textsc{mpmath} library; without it, the check verifies the stored values
(\texttt{results/checks/}).}
\footnotetext[2]{Without a GPU, the check verifies the stored GPU outputs (\texttt{results/checks/}).}
\footnotetext[3]{The replay needs the archive of the runs (\texttt{ledger-archive/}, outside the repository);
without it, the check verifies the invariants of the ledger and reports the replay as skipped.}
\end{table*}

% =================================================================================================================
\section{Data and code availability}
\label{app:data}

The code, the data and the checks are in the repository ReBar, \url{https://github.com/PaoloRibeca/ReBar},
whose release 0.4.3 is archived at Zenodo, \url{https://doi.org/10.5281/zenodo.23750242}; this paper is
arXiv:2610.39307. The
code under the GNU General Public License, version~3 or later, the data and the documentation under the Creative
Commons Attribution 4.0 International licence, and the text of this paper under the Creative Commons
Attribution-NonCommercial-NoDerivatives 4.0 International licence. The repository is the Supplemental Material of this paper; its file
\texttt{SUPPLEMENTAL.md} says where each of the following items is.
\begin{enumerate}
\item \emph{Graph lists}: $\Gamma^{(4)}$ and $\Gamma^{(2)}$ to ten loops, in the notation of
Refs.~\cite{Ribeca2003thesis,GuidaRibeca2006} (\texttt{graphs/lists/}) and, in the form the integrator reads
(\texttt{integrator/data/lists/}), with the canonical Nickel index, the weight $E!/|\mathrm{Aut}\,G|$ at $N=1$ and
the one-vertex-reducible and self-loop flags of every graph; the $O(N)$ weight polynomials are in the ledger, class by
class, with the SHA-256 of every list the integrator ran on. The lists to seven ($\Gamma^{(4)}$) and eight
($\Gamma^{(2)}$) loops coincide with those of Ref.~\cite{Phi4tools2024}, version~1.0.2 (check G3), and the four
six- and seven-loop lists of the integrator are its files, with its columns, reproduced under its MIT licence; the
$\Gamma^{(4)}$ lists from eight loops and the $\Gamma^{(2)}$ lists from nine loops on are new.
\item \emph{Per-class results} (\texttt{ledger/}): the statistics of every class pooled over the runs (number of
samples, mean and central moments to the fourth, extremes, round-off indicators, time per sample), the exact weight
polynomials, the per-order sums with the reference values and the pulls, and a record of every run with its own
statistics (Sec.~\ref{sec:methods-ledger}); and, in the archive deposited with the release (too large for the
repository), the output of every run as received, with the integral, error, samples and diagnostics of every class. With them the coefficients can be recombined for any $N$, the error propagation of
Sec.~\ref{sec:resummation-errors} can be repeated, and two runs can be compared class by class.
\item \emph{Run records}: dated snapshots of the series and tables derived from the ledger
(\texttt{results/coefficients/}), and the reference data and inputs of the checks (\texttt{results/checks/}). The
complete records of the campaigns of Sec.~\ref{sec:results} (reports, logs and environment records) are kept outside
the repository; the run files of the ledger's contributions are deposited with its archive.
\item \emph{Series} (\texttt{series/}): the tools that validate and pool the runs and derive from the ledger the
polynomials $B_L(N)$, $H_L(N)$, $C_L(N)$ and $p_{X,k}(N)$, the coefficients and the generated tables of this paper; the
polynomials $p_{X,k}(N)$ with the full covariance of their coefficients (Sec.~\ref{sec:results-new}), in
\texttt{series/generated/polynomials.json} and, one row per coefficient, \texttt{polynomials.tsv}, with a dated copy
in \texttt{results/coefficients/}.
\item \emph{Checks} (\texttt{checks/}): the scripts of Appendix~\ref{app:reproducibility}.
\item \emph{Code} (\texttt{integrator/}, \texttt{graphs/}): the program that builds and integrates the subtracted integrands (C++17 and CUDA from a single
source; versions 0.3.0 to 0.4.3 for the runs of Sec.~\ref{sec:results}, each recorded with its build signature in
the ledger), with
its self-tests and validation driver, and, in \texttt{results/checks/}, the golden reference files of
Sec.~\ref{sec:methods-kernel}; the graph generator of
Sec.~\ref{sec:methods-generation}; and the scripts that convert the graphs to Nickel indices and evaluate the sum
rules of Appendix~\ref{app:sumrules}.
\end{enumerate}
Every results archive records the hardware, driver and compiler versions and the checksums of all sources. Since the
random numbers of a class depend only on the seed, the kind of integrand, the loop order and the Nickel index
(Sec.~\ref{sec:methods-rng}), any class of any run can be recomputed independently: bit for bit with the same build
on CPUs, for any number of threads, and within round-off with other compilers and libraries and on GPUs
(Sec.~\ref{sec:methods-kernel}).
